%% Authors: Please do not make any changes on the commands below.
%%%%%%%%%%%%%%%%%%%%%%%%%%%%%%%%%%%%%%%%%%%%%%%%%%%%%%%%%%%%%%%%%%

\documentclass[12pt,a4paper]{icicel}

%\textheight=8.2 true in
\setcounter{page}{1}
\topmargin -20pt

\textwidth 16cm
\textheight 25.5cm

\oddsidemargin 0cm
\evensidemargin -0.06cm

\def\currentvolume{0}
\def\currentissue{0}
\def\currentyear{2026}
\def\currentmonth{July}
\def\ppages{1--ICICIC-2026-000}

\newcommand{\BOX}{\hfill $\Box$}
\newcommand{\NAB}{\hfill $\nabla \nabla \nabla$}
\newcommand {\done} {\quad\vrule height4pt WIDTH4PT}
\newcommand{\BYD}{\stackrel{\Delta}{=}}

\newtheorem{corollary}{Corollary}[section]
\newtheorem{theorem}{Theorem}[section]
\newtheorem{lemma}{Lemma}[section]
\newtheorem{proposition}{Proposition}[section]
\newtheorem{remark}{Remark}[section]
\newtheorem{definition}{Definition}[section]
\newtheorem{example}{Example}[section]

\usepackage[dvips]{graphicx}
\usepackage{epsfig}
\usepackage{epsf}

%% Authors: Please do not make any changes on the commands above.
%%%%%%%%%%%%%%%%%%%%%%%%%%%%%%%%%%%%%%%%%%%%%%%%%%%%%%%%%%%%%%%%%%

\title[ICICIC 2026]
      {Positioning EQ-Fuzzy Against Existing Emotion Benchmarks}

\author[N. Shirahama, Y. Yoshimoto, N. Nakaya, J. Nishino and S. Watanabe]{}

\begin{document}

\maketitle

\centerline{\scshape Naruki Shirahama$^1$, Yuma Yoshimoto$^2$, Naofumi Nakaya$^3$,}
\centerline{\scshape Junji Nishino$^4$ and Satoshi Watanabe$^5$}
\medskip
{\footnotesize
\centerline{$^1$Shimonoseki City University}
\centerline{Shimonoseki, Yamaguchi, Japan}
\centerline{nshirahama@ieee.org}}
\medskip
{\footnotesize
\centerline{$^2$NIT (KOSEN), Kitakyushu College}
\centerline{Kitakyushu, Fukuoka, Japan}
\centerline{yoshimoto@kct.ac.jp}}
\medskip
{\footnotesize
\centerline{$^3$Juntendo University}
\centerline{Urayasu, Chiba, Japan}
\centerline{n.nakaya.ac@juntendo.ac.jp}}
\medskip
{\footnotesize
\centerline{$^4$Tamagawa University}
\centerline{Machida, Tokyo, Japan}
\centerline{nishinojunji@eng.tamagawa.ac.jp}}
\medskip
{\footnotesize
\centerline{$^5$Shizuoka Institute of Science and Technology}
\centerline{Fukuroi, Shizuoka, Japan}
\centerline{satoshi\_watanabe@ieee.org}}
\medskip

\centerline{Received Month 2026; accepted Month 2026}

%%\centerline{(Communicated by xxxx)}

\medskip

\begin{abstract}
{\em
Existing emotion and emotional-intelligence benchmarks commonly emphasize
accuracy, agreement, empathy, or classification-style performance.
This paper positions EQ-Fuzzy as a complementary benchmark method for large
language model emotion scoring.
Rather than claiming that EQ-Fuzzy replaces existing benchmarks, the study asks
what additional signals it can expose when benchmark stimuli, scoring rules,
and evaluation targets are compared in a matched setting.
The ICICIC 2026 experiment combines a benchmark coverage matrix with a
stable six-model English EQ-Fuzzy matched subset:
6 models $\times$ 3 literary texts $\times$ 2 evaluation targets $\times$
10 repetitions, yielding 360 response rows and 1440 emotion rows with
valid-output rate 1.0.
The resulting descriptors include valid-output rate, cell means,
within-cell dispersion, fuzzy entropy, profile entropy, and target-shift
magnitude.
An external mini-comparison against curated public EQ-Bench and EmoBench items
is used only as descriptive provenance for conventional benchmark behavior.
The contribution is a compact benchmark-positioning analysis of
uncertainty-aware, stability-oriented, and target-controllable evaluation
signals.
}\\
{\bf Keywords:} Emotion benchmark, Large language model, Fuzzy emotion scoring,
Uncertainty aware evaluation, Benchmark positioning.
\end{abstract}

\section{Introduction}
Large language models are increasingly evaluated on emotional understanding,
empathy, and affective reasoning tasks.
Benchmarks such as EQ-Bench and EmoBench provide useful ways to compare models
on dialogue-based emotional intensity prediction or emotional-intelligence
question answering [1,2].
Psychometric emotional-understanding tests offer another reference point by
comparing model answers with human norm frames [3].
These benchmarks are valuable, but their primary outputs often remain centered
on correctness, agreement, leaderboard score, or human-level comparison.
Such outputs are useful when the research question is whether a model solved a
benchmark item.
They are less direct when the question is how a model distributes affective
intensity over several emotions, how stable that distribution is, or whether
the distribution changes when the requested target of interpretation changes.

EQ-Fuzzy addresses a narrower and complementary question.
For a fixed text and emotion scale, it records continuous emotion intensities
and permits derived descriptors of uncertainty, repeat stability, and target
control.
The ICICIC 2026 workstream therefore asks what EQ-Fuzzy captures beyond
existing emotion benchmarks, not whether it is universally superior to them.
This distinction is important for keeping the paper as a benchmark-positioning
study rather than as a broad leaderboard claim.
It also separates the present paper from adjacent EQ-Fuzzy work on
human-grounded multilingual alignment and persona-temperature effects [4,5].
The paper makes three contributions.
First, it records a reproducible coverage matrix that locates EQ-Fuzzy relative
to EQ-Bench, EmoBench, and SECEU-style emotional-understanding tests.
Second, it reports a stable six-model matched subset in which the same texts,
emotion dimensions, response schema, and repetitions are used across all
models.
Third, it demonstrates reader-side versus character-side target control as a
descriptor that conventional single-target benchmark scores do not naturally
expose.

\section{Benchmark Context}
EQ-Bench evaluates emotional intelligence through dialogue situations and
emotion-intensity judgments.
Its output is a conventional score for a fixed benchmark task [1].
EmoBench evaluates emotional understanding and application with
question-answering items, so its natural output is closer to item accuracy than
to a distributional description of an emotion profile [2].
SECEU and related emotional-understanding studies use psychometric test
frames, making them useful reference points for human-norm comparison but less
directly suited to repeated target-control experiments [3].

These benchmark families are intentionally different from EQ-Fuzzy.
They help answer whether a model gives the expected answer under a predefined
task.
EQ-Fuzzy instead asks how strongly a model scores several emotions for the
same text, how stable those scores are under repeated sampling, and how the
profile changes when the requested evaluation target changes.
The comparison is therefore not a replacement argument.
It is a positioning argument: EQ-Fuzzy contributes descriptors that are not
usually primary outputs of existing emotion benchmarks.
This distinction also shapes the evaluation criteria.
The relevant question is not whether EQ-Fuzzy achieves a higher benchmark
score than the alternatives.
The relevant question is whether it makes additional, reproducible descriptors
available while remaining clear about what those descriptors can and cannot
establish.

\begin{table}[t]
\centering
\footnotesize
\caption{Compact benchmark-positioning matrix generated from the ICICIC comparison artifact.}
\label{tab:icicic-coverage}
{\tiny\setlength{\tabcolsep}{0pt}\begin{tabular}{p{0.28\linewidth}p{0.23\linewidth}p{0.18\linewidth}p{0.23\linewidth}}
\hline
Benchmark & Output & Uncertainty & Target \\
\hline
EQ-Bench & emotional-intensity score & not primary & fixed dialogue task \\
EmoBench & EI question correctness & not primary & fixed EI question \\
SECEU / Emotion Understanding & human-norm score & not primary & fixed psychometric construct \\
EQ-Fuzzy matched subset & descriptor layer & native descriptor & reader-side vs. character-side \\
\hline
\end{tabular}
}
\end{table}

\begin{figure}[t]
\centering
\fbox{\begin{minipage}{0.92\linewidth}
\footnotesize
\begin{tabular}{p{0.42\linewidth}p{0.42\linewidth}}
Conventional benchmark output &
EQ-Fuzzy descriptor layer \\
\hline
correctness, agreement, native task score &
valid-output rate, cell dispersion, fuzzy entropy, profile entropy, target shift
\end{tabular}
\end{minipage}}
\caption{Conceptual role of EQ-Fuzzy as a descriptor layer beside conventional benchmark outputs.}
\label{fig:descriptor-layer}
\end{figure}

\section{Matched Comparison Method}
The comparison is organized around four benchmark families:
EQ-Bench, EmoBench, SECEU / Emotion Understanding, and the EQ-Fuzzy matched
subset.
The coverage matrix compares task type, stimulus type, language, response
format, scoring rule, human reference, repeatability, uncertainty output,
controllability, persona-temperature handling, and reproducibility or license
status.
This matrix is generated from the repository script so that table entries can
be regenerated from the same benchmark-positioning specification.
Table \ref{tab:icicic-coverage} gives a compact manuscript-facing view of this
matrix; the full generated CSV retains the broader audit dimensions.
The matrix is intentionally qualitative.
It records whether a benchmark family exposes a property as a primary output,
not whether that property could be approximated through a secondary analysis.
This conservative coding avoids overstating EQ-Fuzzy's novelty while still
showing where its artifact chain differs.

The EQ-Fuzzy matched subset uses three English literary texts from the shared
registry, four emotions, and a stable six-model panel selected from the
repository model registry.
The stable6 outputs were collected on April 27, 2026 (JST), using OpenRouter
model identifiers recorded in the model registry and generated CSV artifacts.
Manuscript tables use short aliases for layout, while the artifact tables keep
the exact identifiers.
The panel was adopted from the passed SCIS main panel after the earlier
budget4 pilot showed schema noncompliance for Qwen.
Two evaluation targets are compared under the same text and response schema:
\texttt{reader\_side}, which rates emotions felt by a reader after reading the
whole text, and \texttt{character\_side}, which rates emotions plausibly
experienced by the central character or speaker.
The arithmetic is explicit:
6 models $\times$ 3 texts $\times$ 2 target modes $\times$ 10 repetitions
= 360 response rows; multiplying each response by four emotions yields
1440 emotion rows.
A response row is one schema-checked model-text-target-repetition output.
An emotion row is one response row expanded to one emotion dimension.
A model-text-target-emotion cell aggregates 10 repeated scores.
For this run, all 360 response rows were valid.
This design makes target controllability visible without rerunning the
human-grounded multilingual alignment analysis used in the ICECCME workstream.
Each matched-subset response is required to return the same four emotion
scores and short text-grounded reasons.
The reasons are retained for auditability, but the ICICIC analysis centers on
the structured score fields so that the descriptors can be regenerated from
raw JSONL records.

The external mini-comparison uses only manually curated public items whose
source URLs and reuse notes are recorded in the repository configuration.
The current curation includes EQ-Bench v2 and English EmoBench items.
SECEU remains a coverage-matrix comparator because item reuse terms were not
treated as explicitly clear for an external mini-run in this phase.

\section{EQ-Fuzzy Added Descriptors}
Let $x_{m,t,c,r,e}$ denote a 0--100 score produced by model $m$ for text $t$,
target condition $c$, repetition $r$, and emotion $e$.
For each model-text-target-emotion cell, the analysis reports the mean score,
standard deviation, interquartile range, and valid-output rate.
Each score is also mapped to Low, Medium, and High fuzzy membership values, and
a normalized fuzzy entropy descriptor is computed under the locked membership
configuration used by the repository.
The primary \texttt{sigmoid\_s\_v1} membership family uses Low as a smooth
z-function over [0,30], Medium as a smooth pi-function with rise [10,30],
plateau [30,50], and fall [50,70], and High as a smooth s-function over
[50,100].
Fuzzy entropy is $H=-\sum_i \mu_i \log_2 \mu_i$ with $0\log_2 0=0$, normalized
by the dense-grid maximum for the locked family.

The four-emotion profile entropy describes how concentrated or distributed the
emotion profile is within a response.
It is computed by normalizing nonnegative emotion scores as
$p_e=\max(x_e,0)/\sum_e\max(x_e,0)$ and applying
$H_p=-\sum_e p_e\log_2 p_e/\log_2 4$; an all-zero profile is assigned entropy
0.
The reader-versus-character target shift is computed as the absolute difference
between the target-mode cell means for the same model, text, and emotion.
Thus Table \ref{tab:icicic-target-shift} uses $n=3$ paired text cells for each
model-emotion row, whereas Table \ref{tab:icicic-valid-output} uses $n=12$
model-target cells from 3 texts $\times$ 4 emotions.
These descriptors are interpreted descriptively.
They are meant to be read together.
A model can be schema-compliant but unstable, stable but target-insensitive, or
target-sensitive only for particular emotions.
This is why the ICICIC analysis reports several descriptors rather than
collapsing the run into one aggregate score.
They do not prove human validity by themselves and do not replace the
factorial persona-temperature deconfounding analysis reserved for the SCIS
workstream.

\section{Results}
The stable6 matched subset produced the manuscript-facing ICICIC results.
All result tables in this section are generated from the analysis artifacts
rather than manually edited in the manuscript source.
This keeps the paper text separate from the numerical source of record.

Table \ref{tab:icicic-model-descriptors} summarizes model-level descriptors.
The valid-output column shows that the final matched-subset run was suitable
for descriptor analysis across the stable model panel.
The remaining columns show that the models are not interchangeable even when
they all satisfy the response schema.
Within-cell dispersion, fuzzy entropy, profile entropy, and mean target shift
vary by model, which is the signal that a single leaderboard-style score would
hide.
For example, mean target shift is largest for Gemini and Claude, while Grok,
Llama, and GPT show much smaller shifts.
Conversely, Grok has the highest profile entropy despite its small target
shift, and DeepSeek has the largest within-cell dispersion.
The table therefore separates target sensitivity, profile concentration, and
repeat stability instead of collapsing them into a winner-take-all ranking.

\begin{table}[t]
\centering
\footnotesize
\caption{Model-level EQ-Fuzzy descriptors generated from the stable6 matched subset.}
\label{tab:icicic-model-descriptors}
{\tiny\setlength{\tabcolsep}{1pt}\begin{tabular}{llllll}
\hline
model id & valid output rate & mean within cell sd score & mean cell H norm & mean profile entropy & mean target shift abs \\
\hline
anthropic/claude-sonnet-4.5 & 1.0 & 3.197648 & 0.469909 & 0.851784 & 18.716667 \\
deepseek/deepseek-v3.2 & 1.0 & 7.268133 & 0.270318 & 0.774712 & 12.466667 \\
google/gemini-2.5-pro & 1.0 & 5.978447 & 0.269219 & 0.743199 & 24.291667 \\
meta-llama/llama-4-maverick & 1.0 & 4.031645 & 0.37615 & 0.889589 & 5.916667 \\
openai/gpt-5.4 & 1.0 & 2.676327 & 0.365859 & 0.863257 & 6.641667 \\
x-ai/grok-4.20 & 1.0 & 2.403945 & 0.425034 & 0.931072 & 5.75 \\
\hline
\end{tabular}
}
\end{table}

Table \ref{tab:icicic-target-shift} reports the same target-control idea at a
finer resolution.
The reader-side and character-side prompts use the same texts, emotions, schema,
and model panel.
Under this fixed text, schema, model, and repetition protocol, the observed
shift is operationally attributable to the target-mode prompt rather than to a
different benchmark item.
The largest shifts concentrate in particular model-emotion cells, especially
where the reader's post-reading affect and the central character's plausible
affect need not coincide.
Across the table, anger and interest drive many of the larger shifts, whereas
sadness shifts are often smaller.
This pattern is consistent with the task distinction: a reader may feel
interest or surprise toward a scene even when the character or speaker's
plausible affect is anger or sadness.
This is the clearest ICICIC example of the added value of target-controllable
emotion scoring.
It matters for benchmark design because reader-affect and character-affect
prompts may both look like emotion-scoring tasks, but they generate measurably
different score profiles.
EQ-Fuzzy makes that distinction inspectable rather than leaving it implicit
inside an aggregate benchmark score.
The target-shift table also makes the paper's scope concrete.
It does not say that one target mode is correct and the other is wrong.
Instead, it shows that the two targets are experimentally distinguishable and
that their difference can be measured under a fixed text and response schema.

\begin{table}[t]
\centering
\footnotesize
\caption{Target-shift summary by model and emotion generated from the stable6 matched subset.}
\label{tab:icicic-target-shift}
{\tiny\setlength{\tabcolsep}{2pt}\begin{tabular}{lllll}
\hline
Model & Emotion & n & Mean shift & Max shift \\
\hline
Claude & anger & 3 & 27.3 & 57.2 \\
Claude & interest & 3 & 25.667 & 57.0 \\
Claude & sadness & 3 & 7.167 & 10.4 \\
Claude & surprise & 3 & 14.733 & 42.8 \\
DeepSeek & anger & 3 & 24.333 & 47.0 \\
DeepSeek & interest & 3 & 13.833 & 39.0 \\
DeepSeek & sadness & 3 & 1.5 & 2.5 \\
DeepSeek & surprise & 3 & 10.2 & 15.1 \\
Gemini & anger & 3 & 36.833 & 85.5 \\
Gemini & interest & 3 & 27.333 & 43.0 \\
Gemini & sadness & 3 & 8.5 & 15.0 \\
Gemini & surprise & 3 & 24.5 & 48.0 \\
Llama & anger & 3 & 11.0 & 28.0 \\
Llama & interest & 3 & 3.0 & 8.0 \\
Llama & sadness & 3 & 3.667 & 9.0 \\
Llama & surprise & 3 & 6.0 & 9.0 \\
GPT & anger & 3 & 3.667 & 5.6 \\
GPT & interest & 3 & 17.233 & 37.1 \\
GPT & sadness & 3 & 1.6 & 2.7 \\
GPT & surprise & 3 & 4.067 & 7.7 \\
Grok & anger & 3 & 10.833 & 18.5 \\
Grok & interest & 3 & 2.5 & 4.0 \\
Grok & sadness & 3 & 4.667 & 9.0 \\
Grok & surprise & 3 & 5.0 & 9.5 \\
\hline
\end{tabular}
}
\end{table}

Table \ref{tab:icicic-valid-output} gives a direct schema-compliance check by
target mode.
For the accepted stable6 run, the matched-subset descriptors are not driven by
missing or invalid cells.
This matters because an earlier budget4 pilot was archived as provenance after
schema noncompliance appeared in the panel.
The stable6 panel is therefore used for the manuscript-facing result
interpretation.
The archived pilot remains useful because it shows why schema compliance is a
substantive part of benchmark engineering.
However, all manuscript claims about the matched-subset descriptors are based
on the stable6 analysis rather than on the pilot.

\begin{table}[t]
\centering
\footnotesize
\caption{Valid-output rate by model and target mode generated from the stable6 matched subset.}
\label{tab:icicic-valid-output}
{\tiny\setlength{\tabcolsep}{2pt}\begin{tabular}{lllll}
\hline
Model & Target & n & Mean valid & Min valid \\
\hline
Claude & character\_side & 12 & 1.0 & 1.0 \\
Claude & reader\_side & 12 & 1.0 & 1.0 \\
DeepSeek & character\_side & 12 & 1.0 & 1.0 \\
DeepSeek & reader\_side & 12 & 1.0 & 1.0 \\
Gemini & character\_side & 12 & 1.0 & 1.0 \\
Gemini & reader\_side & 12 & 1.0 & 1.0 \\
Llama & character\_side & 12 & 1.0 & 1.0 \\
Llama & reader\_side & 12 & 1.0 & 1.0 \\
GPT & character\_side & 12 & 1.0 & 1.0 \\
GPT & reader\_side & 12 & 1.0 & 1.0 \\
Grok & character\_side & 12 & 1.0 & 1.0 \\
Grok & reader\_side & 12 & 1.0 & 1.0 \\
\hline
\end{tabular}
}
\end{table}

\section{External Mini-Comparison}
The external mini-comparison is included to connect the ICICIC positioning
claim with conventional public benchmark behavior.
It uses curated public EQ-Bench v2 and English EmoBench items whose source
URLs and reuse notes are recorded in the repository.
The analysis reports native EQ-Bench scores where the required score profile
can be parsed, EmoBench exact-match rates, model confidence, and valid-output
rates.
Table \ref{tab:icicic-external-mini} summarizes this provenance-only run.
The external mini-comparison is not used to evaluate model quality.
It is used only to document how public benchmark-style formats behave under
the same artifact discipline.

The external results should be read descriptively.
They are not used to rank the model panel, and they do not override the
matched-subset results.
The table contains 144 response rows, of which 130 are valid.
The invalid rows come from Llama under public-benchmark response formats.
Because the external mini-comparison is provenance rather than the main
experiment, those rows are reported as a limitation instead of being retried in
this phase.
This handling is deliberately conservative.
Rerunning failed external rows could improve a table, but it would blur the
distinction between the accepted matched-subset experiment and the smaller
public-item comparison.
For ICICIC, the important observation is that the same model can be fully
schema-compliant in the EQ-Fuzzy matched subset while being less reliable under
a different public benchmark response format.

\begin{table}[t]
\centering
\footnotesize
\caption{External mini-comparison summary generated from curated public benchmark items; provenance only, not a model-quality ranking.}
\label{tab:icicic-external-mini}
{\tiny\setlength{\tabcolsep}{1pt}\begin{tabular}{p{0.21\linewidth}lll p{0.26\linewidth}p{0.16\linewidth}}
\hline
Source & Rows & Valid & Valid rate & Invalidity & Role \\
\hline
EQ-Bench v2 mini & 72 & 62 & 0.861 & llama-4-maverick format/parse failures & provenance only \\
EmoBench English mini & 72 & 68 & 0.944 & llama-4-maverick format/parse failures & provenance only \\
Total & 144 & 130 & 0.903 & not retried; limitation & descriptive only \\
\hline
\end{tabular}
}
\end{table}

\section{Reproducibility and Artifact Discipline}
The ICICIC artifact chain is designed so that the manuscript can be checked
against stored manifests and generated summaries.
The accepted matched-subset result is the stable6 main run and its stable6
analysis directory.
The corresponding raw JSONL, analysis CSV files, manuscript-facing tables, and
diagnostic plots retained in the artifact directory are all stored under
ICICIC paths.
No result used for the ICICIC paper is written into the ICECCME result tree.
The Qwen-containing budget4 run is retained as archived provenance, but it is
not used for the ICICIC main tables, figures, or claims.
The external mini-comparison is also separated from the final matched-subset
evidence.
Keeping final, pilot, and provenance roles distinct makes the paper auditable
and prevents accidental mixing of result chains.

\section{Implications for Benchmark Design}
The ICICIC results suggest three practical design lessons for emotion
benchmarking with large language models.
First, schema compliance is not a clerical detail.
If a benchmark asks for structured emotional intensities, invalid or
shape-incompatible responses directly affect whether dispersion, entropy, and
target-shift descriptors can be computed.
The stable6 main run therefore matters not only because it produced complete
rows, but because it made the downstream descriptors interpretable.

Second, repeated sampling should be treated as a benchmark design feature when
the target is affective scoring.
Emotion judgments over literary texts are not simply correct or incorrect in
the same way as multiple-choice answers.
A model that returns plausible scores may still vary in a way that changes the
interpretation of its emotional profile.
The within-cell dispersion and fuzzy-entropy descriptors give a compact way to
report this stability question without making a stronger claim about human
validity.

Third, target specification should be explicit.
Reader-side post-reading affect and inferred character-side affect are both
reasonable tasks, but they are not the same task.
When a benchmark does not distinguish them, target ambiguity can be hidden
inside an aggregate score.
The ICICIC matched subset shows that holding text, model, schema, and
repetition count fixed while changing only the target mode creates an auditable
measure of target sensitivity.
This lesson is independent of any particular model ranking and is the main
reason EQ-Fuzzy is useful as a complement to existing benchmarks.

These lessons are intentionally modest.
They do not imply that every emotion benchmark should use literary texts or
fuzzy membership functions.
They imply that benchmark reports should distinguish item correctness,
response validity, repeat stability, uncertainty representation, and target
control when those questions are relevant.
In that sense, EQ-Fuzzy is best understood as a descriptor layer that can sit
beside conventional benchmark scores rather than as a competing leaderboard.

\section{Discussion and Limitations}
This paper is designed as a compact benchmark-positioning study.
The principal finding is that EQ-Fuzzy adds a different kind of evidence to
the evidence produced by existing emotion benchmarks.
It exposes whether repeated responses are stable, whether fuzzy memberships
remain concentrated or ambiguous, and whether a model changes its emotion
profile when the target is switched from reader-side affect to character-side
affect.
These are useful descriptors for benchmark positioning, but they are not a
claim of universal superiority.

The limitations follow directly from the design.
The matched subset is intentionally small, English-only, and based on literary
texts.
It supports a focused comparison of descriptors, not a broad model leaderboard.
The paper therefore should be judged by whether the benchmark-positioning
argument is coherent and reproducible, not by whether it exhausts the space of
emotion benchmarks or literary stimuli.
The external mini-comparison uses only public items with clear reuse notes, so
SECEU remains coverage-matrix-only in this version.
The invalid external mini-comparison rows are retained in the analysis as
valid-output evidence rather than hidden through reruns.
Finally, ICICIC does not establish standalone human validity; that question
belongs to the human-grounded alignment workstream.
Persona-temperature main effects and interactions are also outside the center
of this paper and remain part of the SCIS factorial deconfounding workstream.

\section{Conclusions}
ICICIC 2026 positions EQ-Fuzzy as a complementary method for describing
emotion-evaluation signals that are not always primary outputs of existing
emotion benchmarks.
The contribution is a matched comparison and a reproducible artifact chain
showing how uncertainty-aware, stability-oriented, and target-controllable
descriptors can be reported alongside conventional benchmark positioning.
The results support using EQ-Fuzzy as an additional lens for model emotion
evaluation, while keeping benchmark accuracy, human validity, and
persona-temperature deconfounding as separate questions.
This separation is the main practical value of the paper.
It lets future work combine EQ-Fuzzy descriptors with human-grounded and
factorial analyses without forcing all evaluation goals into a single
benchmark score.

\section*{Acknowledgment}
This work was partially supported by JSPS Grants-in-Aid for Scientific
Research (JSPS KAKENHI), Grant Numbers 26K14983 and 24K13350.

\begin{thebibliography}{10}

\bibitem{paech2023eqbench}
S. J. Paech, EQ-Bench: An emotional intelligence benchmark for large language
models, {\em arXiv preprint arXiv:2312.06281}, 2023.

\bibitem{sabour2024emobench}
S. Sabour, S. Liu, Z. Zhang, J. Liu, J. Zhou, A. Sunaryo, T. Lee,
R. Mihalcea and M. Huang, EmoBench: Evaluating the emotional intelligence of
large language models, {\em Proc. of the 62nd Annual Meeting of the Association
for Computational Linguistics}, pp.5986--6004, 2024.

\bibitem{wang2023ei}
X. Wang, X. Li, Z. Yin, Y. Wu and J. Liu, Emotional intelligence of large
language models, {\em Journal of Pacific Rim Psychology}, vol.17, 2023.

\bibitem{shirahama2025vas}
N. Shirahama, N. Nakaya, S. Watanabe, et al., Assessing emotional intelligence
in AI: Visual Analogue Scale vs. Likert Scale in short story comprehension,
{\em Journal of the Institute of Industrial Applications Engineers}, vol.13,
no.1, pp.21--31, 2025.

\bibitem{shirahama2026math}
N. Shirahama, Y. Yoshimoto, N. Nakaya and S. Watanabe, Mathematical framework
for characterizing emotional individuality in large language models:
Temperature control, fuzzy entropy, and persona-based diversity analysis,
{\em Mathematics}, vol.14, no.7, Art. no.1224, 2026.

\end{thebibliography}

\end{document}
